% © 2026 Ing. David Jaroš — CC BY-NC-ND 4.0
% Licensed under CC BY-NC-ND 4.0. See LICENSE.md.

\ifdefined\INCLUDEMODE
\else
\documentclass[12pt,a4paper]{article}
\usepackage[a4paper,margin=2.5cm]{geometry}
\usepackage{amsmath,amssymb,amsthm}
\usepackage{mathtools}
\usepackage{xcolor}
\usepackage{mdframed}

\newtheorem{theorem}{Theorem}[section]
\newtheorem{proposition}[theorem]{Proposition}
\newtheorem{lemma}[theorem]{Lemma}
\theoremstyle{remark}
\newtheorem{remark}[theorem]{Remark}

\newmdenv[backgroundcolor=green!15,linecolor=green!60!black,linewidth=1.5pt]{resultbox}

\title{GAP-U1 Proof Details}
\author{Ing.\ David Jaroš}
\date{2026-07-14}

\begin{document}
\maketitle
\fi

\section{Symmetry and invariance equations}

Fix
\[
\beta_X(A,B):=\operatorname{Re}\operatorname{Tr}(A X B^\dagger),\qquad X\in M_2(\mathbb C).
\]
Under $A\mapsto UAV$, $B\mapsto UBV$,
\begin{align*}
\beta_X(UAV,UBV)
&=\operatorname{Re}\operatorname{Tr}(UAVX(UBV)^\dagger)\\
&=\operatorname{Re}\operatorname{Tr}(UAVXV^\dagger B^\dagger U^\dagger)\\
&=\operatorname{Re}\operatorname{Tr}(AVXV^\dagger B^\dagger).
\end{align*}
Therefore invariance for all $A,B$ is equivalent to
\[
VXV^\dagger=X\quad\forall V\in K.
\]

Symmetry condition:
\begin{align*}
\beta_X(B,A)
&=\operatorname{Re}\operatorname{Tr}(BXA^\dagger)
=\operatorname{Re}\operatorname{Tr}\big((BXA^\dagger)^\dagger\big)
=\operatorname{Re}\operatorname{Tr}(AX^\dagger B^\dagger).
\end{align*}
Hence $\beta_X(B,A)=\beta_X(A,B)$ for all $A,B$ iff $X=X^\dagger$.

\begin{theorem}[Derivation of the classification criterion]
Invariant symmetric forms in this ansatz are exactly those with
\[
X\in\operatorname{Comm}(K)\cap\operatorname{Herm}(2).
\]
\textbf{Verdict: PROVED.}
\end{theorem}

\section{Case $K=U(2)$}

\begin{proposition}[Commutant for full right unitary group]
If $VXV^\dagger=X$ for all $V\in U(2)$, then $X=\lambda I$.
\textbf{Verdict: PROVED.}
\end{proposition}

\begin{proof}
Take $V=\mathrm{diag}(e^{i\theta_1},e^{i\theta_2})$.
For $X=(x_{ij})$, the relation $VXV^\dagger=X$ gives
\[
(e^{i(\theta_1-\theta_2)}-1)x_{12}=0,\qquad (e^{i(\theta_2-\theta_1)}-1)x_{21}=0
\]
for all $\theta_1,\theta_2$, hence $x_{12}=x_{21}=0$.
Now use
\[
V=\frac{1}{\sqrt2}\begin{pmatrix}1&1\\-1&1\end{pmatrix}
\]
which swaps/rotates basis directions. Invariance then forces $x_{11}=x_{22}$.
So $X=\lambda I$.
\end{proof}

Thus
\[
\beta(A,B)=\lambda\operatorname{Re}\operatorname{Tr}(AB^\dagger).
\]
Uniqueness is only up to real scale.

\begin{theorem}[Uniqueness statement under full symmetry]
For $K=U(2)$ the invariant symmetric bilinear form is unique up to scale.
\textbf{Verdict: PROVED.}
\end{theorem}

\section{Case $K_\eta=U(1)\times U(1)$ with fixed $\eta=\mathrm{diag}(1,-1)$}

Let $V=\mathrm{diag}(e^{i\theta_1},e^{i\theta_2})$.
Condition $VXV^\dagger=X$ gives $x_{12}=x_{21}=0$ and leaves $x_{11},x_{22}$ free.
Hermiticity gives $x_{11},x_{22}\in\mathbb R$.
Write
\[
X=\begin{pmatrix}x_{11}&0\\0&x_{22}\end{pmatrix}
=\alpha I+\beta\eta,
\quad
\alpha=\tfrac{x_{11}+x_{22}}2,
\ \beta=\tfrac{x_{11}-x_{22}}2.
\]

\begin{theorem}[Complete reduced-symmetry family]
For $K_\eta=U(1)\times U(1)$,
\[
\beta_{\alpha,\beta}(A,B)=\operatorname{Re}\operatorname{Tr}\!\big(A(\alpha I+\beta\eta)B^\dagger\big),\qquad \alpha,\beta\in\mathbb R
\]
is the full invariant symmetric family.
\textbf{Verdict: PROVED.}
\end{theorem}

\begin{proposition}[Status of $\operatorname{Re}\operatorname{Tr}(A\eta B^\dagger)$]
The form
\[
\langle A,B\rangle_\eta:=\operatorname{Re}\operatorname{Tr}(A\eta B^\dagger)
\]
is invariant under $K_\eta$ but is not unique there; it is a special point $(\alpha,\beta)=(0,1)$ in a two-parameter family.
It is not invariant under full right $U(2)$.
\textbf{Verdict: PROVED.}
\end{proposition}

\begin{resultbox}
Any uniqueness theorem for GAP-U1 must explicitly include the selected right symmetry group. Without that group choice, uniqueness is false as a universal statement.
\end{resultbox}

\ifdefined\INCLUDEMODE
\else
\end{document}
\fi
